\subsection{Однородные системы линейных уравнений. Теорема о решениях. Теорема о структуре общего решения. Критерий существования нетривиальных решений}

\begin{definition}
Система линейных алгебраических уравнений называется \textbf{однородной}, если столбец свободных членов состоит из нулей ($B = \mathbf{0}$). В матричной форме такая система записывается как:
\[ AX = \mathbf{0} \]
Однородная система всегда совместна, так как всегда имеет \textbf{тривиальное} (нулевое) решение $X = \mathbf{0}$. Все остальные решения (если они существуют) называются \textbf{нетривиальными}.
\end{definition}

\begin{theorem}[О линейной комбинации решений]
Если векторы $X_1, X_2, \dots, X_k$ являются решениями однородной системы $AX = \mathbf{0}$, то любая их линейная комбинация $Y = \alpha_1 X_1 + \alpha_2 X_2 + \dots + \alpha_k X_k$ также является решением этой системы.
\end{theorem}

\begin{tcolorbox}[title=Доказательство, breakable]
Для проверки того, является ли $Y$ решением, подставим его в матричное уравнение системы. Используя свойства дистрибутивности умножения матриц и возможность выноса скаляра за знак произведения, получим:
\[ AY = A(\alpha_1 X_1 + \alpha_2 X_2 + \dots + \alpha_k X_k) = \alpha_1 (AX_1) + \alpha_2 (AX_2) + \dots + \alpha_k (AX_k) \]
Так как по условию $X_i$ — решения однородной системы, то $AX_i = \mathbf{0}$ для всех $i = 1, \dots, k$. Тогда:
\[ AY = \alpha_1 \cdot \mathbf{0} + \alpha_2 \cdot \mathbf{0} + \dots + \alpha_k \cdot \mathbf{0} = \mathbf{0} \]
Следовательно, вектор $Y$ превращает уравнение в верное тождество и является решением.
\end{tcolorbox}

\begin{theorem}[Критерий существования нетривиальных решений]
Однородная система $AX = \mathbf{0}$ имеет нетривиальные решения тогда и только тогда, когда ранг её основной матрицы $R$ меньше числа неизвестных $n$ ($R < n$).
\end{theorem}

В случае квадратной системы ($n \times n$) этот критерий эквивалентен условию $\det A = 0$. Если $R = n$, система имеет только тривиальное решение.

\begin{theorem}[О структуре общего решения]
Пусть дана однородная система уравнений $AX = \mathbf{0}$, ранг матрицы которой $R < n$. Тогда существует набор из $n-R$ линейно независимых решений $X_1, X_2, \dots, X_{n-R}$, называемый фундаментальной системой решений (ФСР), такой, что общее решение системы $X_O$ представляется в виде их произвольной линейной комбинации:
\[ X_O = c_1 X_1 + c_2 X_2 + \dots + c_{n-R} X_{n-R} \]
где $c_i$ — произвольные постоянные.
\end{theorem}

\begin{tcolorbox}[title=Доказательство, breakable]
\textbf{1. Построение ФСР.}
Приведем систему к ступенчатому виду методом Гаусса. Без ограничения общности (путем переименования переменных) будем считать, что первые $R$ переменных $x_1, \dots, x_R$ являются главными, а остальные $n-R$ переменных $x_{R+1}, \dots, x_n$ — свободными. С помощью обратного хода метода Гаусса выразим главные переменные через свободные:
\[
\begin{cases}
x_1 = \beta_{1, R+1}x_{R+1} + \dots + \beta_{1, n}x_n \\
\dots \\
x_R = \beta_{R, R+1}x_{R+1} + \dots + \beta_{R, n}x_n
\end{cases}
\]
Построим $n-R$ частных решений специального вида. Для получения $X_1$ положим свободную переменную $x_{R+1} = 1$, а остальные свободные переменные — нулями. Для $X_2$ положим $x_{R+2} = 1$, остальные свободные — нулями, и так далее.
В результате получим матрицу решений, «хвост» которой (строки с $R+1$ по $n$) образует единичную матрицу $E_{n-R}$.

\textbf{2. Линейная независимость.}
Рассмотрим матрицу $X$, столбцами которой являются векторы $X_1, \dots, X_{n-R}$. В этой матрице содержится минор порядка $n-R$, соответствующий строкам свободных переменных. Этот минор равен определителю единичной матрицы, т.е. 1. Следовательно, ранг матрицы $X$ равен $n-R$, и её столбцы (наши решения) линейно независимы.

\textbf{3. Полнота (всякое решение есть линейная комбинация).}
Пусть $X^* = (\gamma_1, \dots, \gamma_R, \gamma_{R+1}, \dots, \gamma_n)^T$ — произвольное решение системы. Составим вектор $Y = X^* - (\gamma_{R+1} X_1 + \dots + \gamma_n X_{n-R})$. По теореме о линейной комбинации $Y$ также является решением.
Рассмотрим координаты вектора $Y$. Для свободных переменных (индексы от $R+1$ до $n$) по построению векторов $X_i$ получим:
$y_{R+k} = \gamma_{R+k} - \gamma_{R+k} \cdot 1 = 0$.
Таким образом, у решения $Y$ все свободные переменные равны нулю. В методе Гаусса значения главных переменных определяются свободными однозначно. Если все свободные переменные — нули, то и главные переменные обязаны быть нулями. Значит, $Y = \mathbf{0}$, что и доказывает разложение $X^*$ по векторам ФСР.
\end{tcolorbox}

\begin{tcolorbox}[title=Источники, colback=gray!10]
\begin{itemize}
  \item Лекция №\,2, тема «Однородные СЛАУ», временной отрезок 54:00--58:00
  \item Лекция №\,3, тема «Однородные системы. Структура общего решения», временной отрезок 97:00--112:00
  \item PDF «Линейная алгебра\_compressed (1).pdf», раздел 4 «Однородные СЛАУ», стр. 13--15
  \item PDF «Вопросы», вопрос №\,6
\end{itemize}
\end{tcolorbox}
